\section{Discussion}\label{sec:discussion}

\subsection{What fuzzification buys}

The most consistent finding is that representing customers near quintile boundaries by partial memberships improves on hard quintiles, both as raw attributes and as concept lattices, on three datasets and across every parameter setting we tried. The mechanism is simple: about 70\,\% of customers lie near some cut point, and the crisp context separates customers who behave alike while merging customers who do not. The gain is concentrated where it matters commercially, in the top decile of spenders. Robustness and stability point the same way: fuzzy RFM-FCA gives the most stable segment structure of all representations compared, including the crisp lattice and every clustering baseline at comparable granularity.

\subsection{What it does not buy}

Three limits are equally consistent. First, segments are a compression: a plain log-RFM model is at least as accurate for prediction and targeting, so fuzzy RFM-FCA should not be sold as a forecasting engine. Second, most of the predictive gain comes from fuzzifying the context rather than from the lattice; the lattice added spend signal beyond the fuzzy memberships on two datasets and lost it on the third. Third, the full fuzzy lattice has about four times as many segments as the crisp one, and its rules are not more confident on average. Against fuzzy C-means, the natural fuzzy alternative, accuracy is roughly level.

\subsection{Practical recommendation}

These results suggest a two-level use. The full fuzzy lattice is the modelling representation: it predicts and targets better than crisp RFM-FCA and K-means and is the most stable. The 10-concept summary is the communication layer: ten named segments, each defined by about two conditions such as ``fully frequent and high-spending, at least partly very recent'', with the greedy guarantee of Proposition~\ref{prop:greedy}. If only predictions from unnamed groups are needed, fuzzy C-means or log-RFM are simpler and as good. What fuzzy RFM-FCA offers that neither does is segments that are exactly described, ordered, linked by rules and aware of boundary customers.

\subsection{Limitations}

\begin{itemize}
\item Effect sizes over crisp RFM-FCA are small (0.004--0.010 AUC, 0.009--0.012 $R^2$), and on smaller datasets the fuzzy lattice needs tuned regularisation.
\item Interpretability was assessed with structural proxies only; a study with marketers is needed to confirm that the segments are more useful in practice.
\item The method was developed on Online Retail II and confirmed on two further datasets; the later analyses reuse the same three datasets, and the tie-aware membership and the hybrid analysis are exploratory.
\item Only one family of memberships (triangular on the rank scale) and a three-level truth scale were studied; a five-level scale was computationally infeasible with our closure algorithm.
\item All datasets are retail; other sectors and contractual settings remain untested.
\end{itemize}
